%! Author = k
%! Date = 2024/11/20

% Preamble
\documentclass[11pt]{article}

\usepackage{amsmath}
\usepackage[UTF8]{ctex}
\usepackage{multirow}
\usepackage{booktabs} % 用于更好看的表格线
\usepackage{array} % 用于调整列宽
\usepackage{caption} % 用于设置表格标题
\usepackage{geometry}

\pagestyle{empty} % 隐藏页眉页脚

\captionsetup[table]{labelformat=empty} % 禁用表格编号

\geometry{a4paper,paperwidth = 30cm, paperheight = 30cm}

\begin{document}

    % 表格代码开始
    \begin{table}[h]
        \centering
        \renewcommand{\arraystretch}{1.5} % 调整行间距
        \setlength{\tabcolsep}{12pt} % 调整列间距
        \begin{tabular}{
            *{3}{>{\centering\arraybackslash}p{1.5cm}} | *{3}{>{\centering\arraybackslash}p{1.5cm}} % 列格式
        }
            \hline % 表格顶部边框
            \multicolumn{3}{c|}{现$~~~~~~~~~~$态} &
            \multicolumn{3}{c}{次$~~~~~~~~~~$态} \\ % 合并单元格
            \hline
            $Q_2^n$ & $Q_1^n$ & $Q_0^n$ & $Q_2^{n+1}$ & $Q_1^{n+1}$ & $Q_0^{n+1}$ \\ % 表头第二行
            \hline
            % 添加数据行
            1 & 0 & 1 & 0 & 1 & 0 \\
            1 & 1 & 0 & 0 & 1 & 0 \\
            1 & 1 & 1 & 0 & 0 & 0 \\
            \hline
        \end{tabular}
        \caption{表$~~~~~~$无效状态表} % 表格标题
        \label{tab:D_flip_flop} % 表格引用标签
    \end{table}
    % 表格代码结束

\end{document}